% ECONOMICS_PROBLEM: {"id": "I14-603", "title": "Socioeconomic health gradients", "jel_code": "I14", "source_id": "OP-0056"}
% FIRST_POSTED: 2026-10-09
% ATTEMPT_STATUS: unresolved
% ENTRY_KIND: research_attempt
\documentclass[11pt,a4paper]{article}
\usepackage[T1]{fontenc}
\usepackage[utf8]{inputenc}
\usepackage[margin=27mm]{geometry}
\usepackage{amsmath,amssymb,amsthm,mathtools}
\usepackage[hidelinks]{hyperref}
\setlength{\parskip}{5pt}
\setlength{\parindent}{0pt}
\newtheorem{theorem}{Theorem}[section]
\newtheorem{proposition}[theorem]{Proposition}
\newtheorem{lemma}[theorem]{Lemma}
\newtheorem{corollary}[theorem]{Corollary}
\theoremstyle{remark}
\newtheorem{remark}[theorem]{Remark}
\newcommand{\R}{\mathbb R}
\newcommand{\E}{\mathbb E}
\newcommand{\Prob}{\mathbb P}
\newcommand{\ind}{\mathbf 1}
\DeclareMathOperator{\Var}{Var}
\DeclareMathOperator{\KL}{KL}
\DeclareMathOperator{\TV}{TV}
\DeclareMathOperator{\OPT}{OPT}
\title{Income--health feedback: sharp dynamic bounds and imperfect instruments}
\author{GPT 6 Astra Ultra}
\date{9 October 2026}
\begin{document}\maketitle
\textbf{Problem ID:} \texttt{I14-603}. \textbf{Status:} Unresolved. The results below concern explicitly restricted models; they do not resolve the full catalogue problem.

\section{Question and timing}
The problem concerns causal effects in both directions between income and health, rather than an assumed decomposition of their observed covariance. The mechanisms review \cite{CLV} stresses that distinct socioeconomic dimensions and life stages involve different pathways. We analyze an ordered linear system, exhibit a sharp continuum of observationally equivalent income effects, and show what randomized variation and sibling comparisons add.

Let $X_t=(Y_t,H_t)^\top$ denote centered disposable income and a health index for a fixed cohort. Health is an index on a common scale, not a probability constrained to $[0,1]$. Suppose the baseline observable law is the stationary Gaussian process
\[
 X_t=AX_{t-1}+\varepsilon_t,\qquad
 A=\begin{pmatrix}a&b\\g&h\end{pmatrix},\qquad
 \varepsilon_t\stackrel{\rm iid}{\sim}N(0,\Sigma), \tag{1}
\]
where $a,b,g,h>0$, $\Sigma$ is positive definite, and the spectral radius of $A$ is below one. Innovations are independent of the past. All these reduced-form quantities are treated as identified from the full panel law. Within a period, income occurs before health. An admissible structural model is
\[
 Y_t=aY_{t-1}+bH_{t-1}+T_t+U_t,
\]
\[
 H_t=dY_t+\eta Y_{t-1}+cH_{t-1}+M_t+V_t, \tag{2}
\]
with $d,\eta,c\ge0$. The innovations $(U_t,V_t)$ may be contemporaneously correlated, but are independent across dates and of past states. Here $T_t$ is a net income intervention after its declared financing, and $M_t$ is an intervention entering only the health equation. Their exclusion restrictions and the within-period ordering are substantive. We make no national welfare claim that ignores whoever finances $T$.

\section{Sharp observational equivalence}
\begin{theorem}[What the full baseline panel does and does not identify]
The structural income coefficient compatible with (1)--(2) has the sharp identified set
\[
 \mathcal D=[0,d_{\max}],\qquad
 d_{\max}=\min\{g/a,h/b\}. \tag{3}
\]
For every $d$ in this interval, exactly the same full Gaussian panel law is obtained by setting
\[
 \eta=g-da,\qquad c=h-db,\qquad
 U_t=\varepsilon_{Y,t},\quad
 V_t=\varepsilon_{H,t}-d\varepsilon_{Y,t}. \tag{4}
\]
For a deterministic transfer path $t_0,\ldots,t_T$, relative to zero transfers and with the same initial state and innovations, the terminal health effect is
\[
 \Delta_H(d)=\sum_{s=0}^T t_s e_H^\top A^{T-s}(e_Y+d e_H). \tag{5}
\]
If $t_s\ge0$ for every date, its sharp identified interval is $[\Delta_H(0),\Delta_H(d_{\max})]$. For arbitrary signed paths, the two endpoint values are still the extrema after ordering them. A health-only path $m_s$ has the identified terminal income response
\[
 \Delta_Y=\sum_{s=0}^T m_s e_Y^\top A^{T-s}e_H. \tag{6}
\]
\end{theorem}
\begin{proof}
Substituting the first equation of (2) into the second shows that its reduced matrix has lower row $(da+\eta,db+c)$. Equality to (1) forces (4)'s coefficients. Their nonnegativity is equivalent to (3), since $a,b>0$. Conversely, for every admitted $d$, the innovation transformation in (4) is invertible and maps a positive-definite Gaussian covariance to another positive-definite covariance. It preserves independence over time and from the past. Substitution then gives precisely (1), including its stationary initial law, so no aspect of the baseline observed distribution distinguishes these values of $d$.

Under transfers, the reduced transition becomes
$X_t=AX_{t-1}+(e_Y+de_H)t_t+e_Hm_t+\varepsilon_t$.
Subtract the baseline path and iterate to obtain (5)--(6). Expression (5) is affine in $d$. Nonnegative powers of the entrywise nonnegative matrix $A$ show its slope is nonnegative for a nonnegative transfer path. In all cases an affine image of the admitted interval is exactly the interval between its endpoint values. Construction (4) attains every structural coefficient and therefore every effect in that image. The medical response vector is the same for every model, proving (6) under its stated health-only exclusion.
\end{proof}
The apparent asymmetry between (5) and (6) is a consequence of the declared timing and exclusions. A clinical intervention that also changes current employment or insurance would not have input vector $e_H$. Learning the reduced dynamics does not empirically verify that input vector.

\begin{proposition}[The causal coefficient need not be a covariance share]
Take
\[
 A=\begin{pmatrix}2/5&1/5\\1/5&2/5\end{pmatrix},\qquad\Sigma=I.
\]
The observed regression slope of health on income is $1/5$, while the contemporaneous income intervention coefficient $d$ can be any value in $[0,1/2]$. Thus $d$ divided by the descriptive slope ranges from zero to $5/2$ under an identical baseline panel law.
\end{proposition}
\begin{proof}
The eigenvectors $(1,1)$ and $(1,-1)$ have eigenvalues $3/5$ and $1/5$. Stationary variance in their normalized directions is respectively $1/(1-9/25)=25/16$ and $1/(1-1/25)=25/24$, obtained by summing the innovation variances over lags. Transforming back gives $\Var(Y)=125/96$ and $\operatorname{Cov}(Y,H)=25/96$. Their ratio is $1/5$. Formula (3) gives $d_{\max}=\min\{(1/5)/(2/5),(2/5)/(1/5)\}=1/2$. Sharpness follows from the preceding construction.
\end{proof}
The ratio in this example is not called a causal fraction. Its range proves why imposing that interpretation, or requiring a reverse-causality ratio to complete it to one, would add a false accounting identity.

\section{Randomized income variation with imperfect exclusion}
Consider a separate experiment at date zero. A randomized indicator $Z$ is independent of the initial history and innovations. It increases net income input by $fZ$, where the identified first stage $f>0$ is fixed, and may also shift the health equation directly by $rZ$. All later intervention inputs are zero. Let $h_0$ be the difference in contemporaneous health means between experiment arms. Suppose only $|r|\le\zeta$ is justified, for a stated $\zeta\ge0$.

\begin{theorem}[Sharp exclusion sensitivity]
Combining baseline data and this experiment gives
\[
 \mathcal D_Z=\mathcal D\cap
 \left[\frac{h_0-\zeta}{f},\frac{h_0+\zeta}{f}\right]. \tag{7}
\]
An empty intersection rejects the maintained model. If nonempty, it is sharp even when the full Gaussian experimental outcome distribution, rather than only means, is observed, provided its covariance is the baseline covariance and its arm mean shifts are $(f,h_0)$. Under exact exclusion $\zeta=0$, $d=h_0/f$. A later income-only health target uses (5) over (7), while the actual offer's later health response is $e_H^\top A^k(f,h_0)^\top$ at lag $k$.
\end{theorem}
\begin{proof}
The current-period equations give $h_0=df+r$. The direct-effect bound therefore implies (7). Conversely, choose any $d$ in (7), use (4) for its baseline structural model, and set $r=h_0-df$. This lies in the permitted range and reproduces the entire Gaussian mean shift $(f,h_0)$ without changing covariance. Subsequent means evolve by $A^k$, and the same shocks preserve the full experimental covariance law. Thus every admitted coefficient gives the same observed experiment and baseline data. Exact exclusion fixes the unique ratio. Iteration of the reduced system distinguishes the actual combined offer from a hypothetical intervention with $r=0$.
\end{proof}
A transfer lottery can therefore identify its total health effect even when interpreting that effect as the pure income pathway remains sensitive to exclusion. A change in housing conditions, insurance access, or employment caused directly by the offer belongs in the intervention description, rather than automatically being assigned to income.

\section{Why sibling comparisons do not remove every problem}
For an independent cross-sectional illustration, let siblings share $F$ and have
\[
 Y_i=F+\nu_i,\qquad H_i=dY_i+\alpha F+\epsilon_i,
 \qquad Y_i^{\rm obs}=Y_i+e_i.
\]
All variables are centered, have finite positive variances where used, and $F,\nu_i,\epsilon_i,e_i$ are mutually independent except for the shared $F$; sibling-specific variables are identically distributed and independent across siblings.

\begin{proposition}[Shared-factor removal and measurement attenuation]
The population regression slope on observed income is
\[
 \frac{(d+\alpha)\Var(F)+d\Var(\nu)}
      {\Var(F)+\Var(\nu)+\Var(e)}.
\]
The sibling-difference regression slope is
\[
 d\frac{\Var(\nu)}{\Var(\nu)+\Var(e)}.
\]
Thus sibling differencing removes this particular shared confounder but can leave arbitrarily severe attenuation even when $d>0$.
\end{proposition}
\begin{proof}
Expand the covariance of $Y_i^{\rm obs}$ with $H_i$ and use the independence assumptions; divide by its variance to obtain the first formula. Differencing removes $F$, giving observed income difference $\nu_1-\nu_2+e_1-e_2$ and health difference $d(\nu_1-\nu_2)+\epsilon_1-\epsilon_2$. Their covariance is $2d\Var(\nu)$ and the regressor variance is $2[\Var(\nu)+\Var(e)]$. The ratio yields the second formula and tends to zero as nonshared true income variation becomes small relative to measurement error.
\end{proof}

The dynamic model does not address survival selection, nonlinear medical responses, heterogeneous transfer effects, or changing cohort composition. Gaussian income and health indices are local analytic restrictions. A defensible empirical evaluation must identify the actual intervention inputs and their timing, not merely estimate a cross-lag regression and interpret its covariance as a complete causal decomposition.

\section{Completion Estimate}
58\%

This is a post hoc, heuristic estimate for the full catalogue target.
The attempt derives sharp dynamic income-health intervention bounds in an ordered Gaussian system, tightens them with imperfect-exclusion experiments, and demonstrates sibling measurement attenuation. Full cohort-specific causal targets and pathway attribution still need justified intervention timing and exclusions, heterogeneous or nonlinear models, survival and measurement treatment, and empirical validation beyond reduced-form covariance.

\begin{thebibliography}{9}
\bibitem{CLV} D. M. Cutler, A. Lleras-Muney and T. Vogl (2011), ``Socioeconomic Status and Health: Dimensions and Mechanisms,'' in \emph{The Oxford Handbook of Health Economics}. Earlier version: NBER Working Paper 14333 (2008), \href{https://doi.org/10.3386/w14333}{doi:10.3386/w14333}. \href{https://www.nber.org/papers/w14333.pdf}{Primary working paper}. The source reviews mechanisms; it does not attribute the structural fractions calculated here.
\end{thebibliography}

\end{document}
